\documentclass[11pt,a4paper]{article}

\usepackage[utf8]{inputenc}
\usepackage[T1]{fontenc}
\usepackage{amsmath, amssymb}
\usepackage{geometry}
\usepackage{xcolor}
\usepackage{hyperref}
\usepackage{booktabs}
\usepackage{mdframed}
\usepackage{listings}

\geometry{margin=2.2cm}

\hypersetup{colorlinks=true, linkcolor=blue!60!black, urlcolor=blue!60!black}

\definecolor{codebg}{RGB}{247,247,247}
\definecolor{codekw}{RGB}{0,90,180}
\definecolor{codestr}{RGB}{170,60,0}
\definecolor{codecom}{RGB}{100,100,100}

\lstdefinestyle{py}{
    language=Python,
    backgroundcolor=\color{codebg},
    basicstyle=\ttfamily\small,
    keywordstyle=\color{codekw}\bfseries,
    commentstyle=\color{codecom}\itshape,
    stringstyle=\color{codestr},
    numbers=left,
    numberstyle=\tiny\color{gray},
    stepnumber=1,
    numbersep=8pt,
    showstringspaces=false,
    breaklines=true,
    frame=single,
    rulecolor=\color{gray!40},
    tabsize=4,
    xleftmargin=1.2em,
}
\lstset{style=py}

\newcommand{\R}{\mathbb{R}}

\title{\textbf{JAX / Flax / Optax --- Aide-mémoire pour construire un MLP}\\
\large Fonctions utiles et exemples d'implémentation rapide, appliqués au projet Modelization\_V1}
\author{Institut de la Vision --- Directed Evolution Loop}
\date{}

\begin{document}
\maketitle
\tableofcontents
\newpage

% ═══════════════════════════════════════════════════════════════════════════
\section{Introduction}

Ce document rassemble, avec des exemples testés et fonctionnels, les briques
JAX / Flax / Optax nécessaires pour construire, entraîner et régulariser un
MLP --- dans notre cas, pour recommencer la reconstruction des poids du modèle
profile ($F$) puis du modèle de Potts ($F, J$) à partir des cibles de
log-enrichissement déjà construites par \texttt{RegressionV1.build\_multi\_round\_dataset}.

Toutes les versions de bibliothèques utilisées pour tester les exemples de ce
document :
\begin{center}
\begin{tabular}{ll}
\toprule
Bibliothèque & Version testée \\
\midrule
\texttt{jax}   & 0.6.2 \\
\texttt{optax} & 0.2.8 \\
\texttt{flax}  & 0.10.7 \\
\bottomrule
\end{tabular}
\end{center}

\begin{mdframed}
\textbf{Pourquoi Flax + Optax et pas PyTorch ?} Tout le code du projet
(\texttt{sequence\_classesV1.py}, \texttt{RegressionV1.py}) est déjà en JAX.
Rester dans cet écosystème permet de garder un seul système de clés
aléatoires cohérent, de \texttt{jit}-compiler la boucle d'entraînement
complète, et d'éviter d'introduire une seconde bibliothèque de deep learning
pour une seule tâche.
\end{mdframed}

% ═══════════════════════════════════════════════════════════════════════════
\section{Rappels JAX}

\subsection{\texttt{jax.numpy} vs \texttt{numpy}}
\texttt{jax.numpy} (\texttt{jnp}) reproduit l'API de NumPy mais toutes les
opérations sont \textbf{immutables} (pas d'affectation \texttt{arr[i]=x} ;
utiliser \texttt{arr.at[i].set(x)}) et \textbf{différentiables}. C'est la
brique de base sur laquelle tout le reste (Flax, Optax) repose.

\begin{lstlisting}
import jax.numpy as jnp

a = jnp.array([1.0, 2.0, 3.0])
b = a.at[0].set(9.0)          # b = [9,2,3] ; a est inchange (immutable)
c = a + b                     # operations vectorisees comme en NumPy
\end{lstlisting}

\subsection{Clés aléatoires (\texttt{jax.random})}
JAX n'a pas d'état aléatoire global : chaque tirage nécessite une clé
explicite, et une clé ne doit \textbf{jamais être réutilisée} pour deux
tirages différents --- on la \texttt{split} à chaque fois.

\begin{lstlisting}
key = jax.random.key(0)
key, k1, k2 = jax.random.split(key, 3)   # 2 sous-cles independantes + la clef "suivante"
x = jax.random.normal(k1, shape=(10,))
y = jax.random.uniform(k2, shape=(10,))
\end{lstlisting}

\subsection{PyTrees}
Un \emph{pytree} est une structure imbriquée (dict, liste, tuple) de tableaux.
Les paramètres d'un modèle (Flax ou fait main) sont toujours des pytrees.
\texttt{jax.tree\_util.tree\_map} applique une fonction à chaque feuille.

\begin{lstlisting}
params = {'W': jnp.ones((3,3)), 'b': jnp.zeros(3)}
shapes = jax.tree_util.tree_map(lambda p: p.shape, params)
# {'W': (3, 3), 'b': (3,)}
zeroed = jax.tree_util.tree_map(lambda p: p * 0.0, params)
\end{lstlisting}

\subsection{\texttt{jax.grad} et \texttt{jax.value\_and\_grad}}
Calcule le gradient d'une fonction scalaire par rapport à son premier
argument (par défaut). \texttt{value\_and\_grad} renvoie la valeur \emph{et}
le gradient en un seul passage --- toujours préférable dans une boucle
d'entraînement pour éviter de calculer deux fois la passe avant.

\begin{lstlisting}
def f(x):
    return jnp.sum(x ** 2)

g = jax.grad(f)
print(g(jnp.array([1., 2., 3.])))          # [2. 4. 6.]

val, grad = jax.value_and_grad(f)(jnp.array([1., 2., 3.]))
print(val, grad)                           # 14.0  [2. 4. 6.]
\end{lstlisting}

\subsection{\texttt{jax.jit}}
Compile une fonction en XLA à la première exécution (pour chaque
combinaison de formes d'entrée) ; les exécutions suivantes sont beaucoup plus
rapides. \textbf{Toute la boucle d'entraînement} (un pas de descente de
gradient) devrait être compilée en un seul \texttt{jit}, pas les sous-étapes
séparément.

\begin{lstlisting}
@jax.jit
def train_step(params, opt_state, x, y):
    ...
    return params, opt_state, loss
\end{lstlisting}

\subsection{\texttt{jax.vmap}}
Vectorise une fonction écrite pour \emph{une seule} entrée, pour l'appliquer
à un batch, sans boucle Python explicite ni réécriture des indices.

\begin{lstlisting}
def score_one(seq, F):
    return jnp.sum(F[seq, jnp.arange(seq.shape[0])])

batched_score = jax.vmap(score_one, in_axes=(0, None))   # vectorise sur seq, pas sur F
scores = batched_score(seqs, F)   # seqs: (N, L) -> scores: (N,)
\end{lstlisting}

% ═══════════════════════════════════════════════════════════════════════════
\section{Construire une couche dense ``à la main''}

Utile pour comprendre exactement ce que fait Flax en coulisses, ou si l'on
préfère na pas dépendre de Flax du tout.

\begin{lstlisting}
def init_dense(key, in_dim, out_dim):
    """Initialisation de Glorot (Xavier) uniforme."""
    k_w, k_b = jax.random.split(key)
    limit = jnp.sqrt(6.0 / (in_dim + out_dim))
    W = jax.random.uniform(k_w, (in_dim, out_dim), minval=-limit, maxval=limit)
    b = jnp.zeros(out_dim)
    return {'W': W, 'b': b}

def init_mlp(key, layer_sizes):
    """layer_sizes = [140, 64, 32, 1] par exemple."""
    keys = jax.random.split(key, len(layer_sizes) - 1)
    return [
        init_dense(k, d_in, d_out)
        for k, d_in, d_out in zip(keys, layer_sizes[:-1], layer_sizes[1:])
    ]

def mlp_forward(params, x):
    for layer in params[:-1]:
        x = jax.nn.relu(x @ layer['W'] + layer['b'])
    last = params[-1]
    return (x @ last['W'] + last['b']).squeeze(-1)
\end{lstlisting}

\begin{mdframed}
\textbf{Testé} : \texttt{init\_mlp(key, [140, 64, 32, 1])} puis
\texttt{mlp\_forward(params, x)} sur \texttt{x} de forme \texttt{(5, 140)}
produit une sortie de forme \texttt{(5,)} --- confirmé dans ce
document avant rédaction.
\end{mdframed}

% ═══════════════════════════════════════════════════════════════════════════
\section{Fonctions d'activation (\texttt{jax.nn})}

\begin{center}
\begin{tabular}{lll}
\toprule
Fonction & Formule & Usage typique \\
\midrule
\texttt{jax.nn.relu}       & $\max(0,x)$                         & Défaut robuste, rapide \\
\texttt{jax.nn.gelu}       & $x\Phi(x)$ (lissée)                 & Souvent légèrement meilleur que ReLU \\
\texttt{jax.nn.tanh}       & $\tanh(x)$                           & Sorties bornées $[-1,1]$ \\
\texttt{jax.nn.sigmoid}    & $1/(1+e^{-x})$                       & Probabilités, portes \\
\texttt{jax.nn.softplus}   & $\log(1+e^x)$                        & Version lisse de ReLU \\
\texttt{jax.nn.leaky\_relu}& $x$ si $x>0$, sinon $\alpha x$        & Évite les neurones ``morts'' \\
\texttt{jax.nn.softmax}    & $e^{x_i}/\sum_j e^{x_j}$              & Couche de sortie multi-classes \\
\bottomrule
\end{tabular}
\end{center}

\begin{lstlisting}
h = jax.nn.relu(x @ W + b)
p = jax.nn.softmax(logits, axis=-1)
\end{lstlisting}

% ═══════════════════════════════════════════════════════════════════════════
\section{Fonctions de perte}

\subsection{MSE (log-ratio cible)}
\begin{lstlisting}
def mse_loss(preds, targets):
    return jnp.mean((preds - targets) ** 2)
\end{lstlisting}

\subsection{Perte de Poisson (log-vraisemblance négative)}
Utile si l'on entraîne directement sur les comptages bruts plutôt que sur le
log-ratio (cf. \texttt{Poisson\_regression\_AAV.tex}) :
\[
    \ell(\mu, y) = \mu - y\log\mu \qquad (\text{à une constante près}).
\]
\begin{lstlisting}
def poisson_nll(log_mu, y, eps=1e-6):
    """log_mu : sortie brute du reseau (avant exp) ; y : comptage observe."""
    mu = jnp.exp(log_mu)
    return jnp.mean(mu - y * jnp.log(mu + eps))
\end{lstlisting}

\subsection{Huber (robuste aux valeurs aberrantes)}
Combine MSE près de zéro et une pénalité linéaire (moins sensible aux
grosses erreurs isolées) au-delà d'un seuil $\delta$ :
\begin{lstlisting}
def huber_loss(preds, targets, delta=1.0):
    err = preds - targets
    abs_err = jnp.abs(err)
    quad = jnp.minimum(abs_err, delta)
    lin = abs_err - quad
    return jnp.mean(0.5 * quad**2 + delta * lin)
\end{lstlisting}

\subsection{Terme de régularisation L2 manuel}
Si l'on n'utilise pas le \texttt{weight\_decay} d'Optax (\S\ref{sec:optax}),
on peut l'ajouter directement dans la perte :
\begin{lstlisting}
def l2_penalty(params, lam=1e-4):
    leaves = jax.tree_util.tree_leaves(params)
    return lam * sum(jnp.sum(p ** 2) for p in leaves)

def loss_with_l2(params, x, y):
    preds = mlp_forward(params, x)
    return mse_loss(preds, y) + l2_penalty(params)
\end{lstlisting}

% ═══════════════════════════════════════════════════════════════════════════
\section{Flax (\texttt{flax.linen}) : définir un MLP proprement}

Flax évite d'écrire soi-même l'initialisation et l'empilement des couches.
\texttt{nn.compact} permet de déclarer les couches directement dans
\texttt{\_\_call\_\_}.

\begin{lstlisting}
import flax.linen as nn

class ProfileMLP(nn.Module):
    hidden_dim: int = 64
    dropout_rate: float = 0.1

    @nn.compact
    def __call__(self, x, training: bool):
        x = nn.Dense(self.hidden_dim)(x)
        x = nn.relu(x)
        x = nn.Dropout(rate=self.dropout_rate, deterministic=not training)(x)
        x = nn.Dense(1)(x)
        return x.squeeze(-1)
\end{lstlisting}

\subsection{Initialisation des paramètres}
Flax sépare la \emph{définition} du modèle de ses \emph{paramètres} : on
initialise en appelant \texttt{model.init(...)} avec une entrée factice de la
bonne forme.

\begin{lstlisting}
key = jax.random.key(0)
key, k_init, k_drop = jax.random.split(key, 3)

model = ProfileMLP(hidden_dim=32)
dummy_x = jnp.zeros((1, 140))   # (1, L*A) pour le modele profile
variables = model.init({'params': k_init, 'dropout': k_drop}, dummy_x, training=False)
params = variables['params']
\end{lstlisting}

\begin{mdframed}
\textbf{Testé} : \texttt{model.init(...)} sur \texttt{ProfileMLP(hidden\_dim=32)}
avec une entrée \texttt{(1, 140)} produit
\texttt{\{'Dense\_0': \{'bias': (32,), 'kernel': (140, 32)\}, 'Dense\_1': \{'bias': (1,), 'kernel': (32, 1)\}\}}
--- confirmé.
\end{mdframed}

\subsection{Passe avant (\texttt{apply})}
\begin{lstlisting}
preds = model.apply(
    {'params': params}, x,
    training=True,
    rngs={'dropout': k_drop},   # necessaire seulement si training=True (dropout actif)
)
\end{lstlisting}

% ═══════════════════════════════════════════════════════════════════════════
\section{Optax : optimiseurs, schedules, clipping}
\label{sec:optax}

\subsection{AdamW}
\texttt{weight\_decay} dans \texttt{optax.adamw} est \textbf{découplé} du
gradient (contrairement à un $\ell_2$ ajouté dans la perte, qui interagit
avec les taux d'apprentissage adaptatifs d'Adam) --- c'est la raison pour
laquelle AdamW généralise mieux qu'Adam + $\ell_2$-dans-la-perte.

\begin{lstlisting}
import optax

tx = optax.adamw(learning_rate=1e-3, weight_decay=1e-4)
opt_state = tx.init(params)
\end{lstlisting}

\subsection{Schedules de taux d'apprentissage}
\begin{lstlisting}
# Decroissance exponentielle
sched1 = optax.exponential_decay(init_value=1e-3, transition_steps=100, decay_rate=0.9)

# Warmup lineaire puis decroissance cosinus (tres courant)
sched2 = optax.warmup_cosine_decay_schedule(
    init_value=0.0, peak_value=1e-3, warmup_steps=20, decay_steps=300,
)

tx = optax.adamw(learning_rate=sched2, weight_decay=1e-4)
\end{lstlisting}

\begin{mdframed}
\textbf{Testé} : \texttt{warmup\_cosine\_decay\_schedule} avec
\texttt{warmup\_steps=20, decay\_steps=200} donne
$\eta(0){=}0$, $\eta(20){\approx}10^{-3}$ (pic), puis décroît vers $0$ à
l'étape $200$ --- confirmé.
\end{mdframed}

\subsection{Gradient clipping (stabilité)}
Utile si les gradients explosent (rare pour un petit MLP, mais bon réflexe).
Se compose avec l'optimiseur via \texttt{optax.chain} :
\begin{lstlisting}
tx = optax.chain(
    optax.clip_by_global_norm(1.0),
    optax.adamw(learning_rate=1e-3, weight_decay=1e-4),
)
\end{lstlisting}

\subsection{Weight decay sélectif (exclure les biais)}
Il est courant de ne pas régulariser les biais (uniquement les poids). Avec
Flax, on peut construire un masque miroir de l'arborescence des paramètres :
\begin{lstlisting}
import flax.traverse_util as traverse_util

def build_mask(params):
    flat = traverse_util.flatten_dict(params)
    mask_flat = {k: (k[-1] != 'bias') for k in flat}
    return traverse_util.unflatten_dict(mask_flat)

mask = build_mask(params)
tx = optax.chain(
    optax.adamw(learning_rate=1e-3, weight_decay=0.0),   # weight decay applique separement
    optax.masked(optax.add_decayed_weights(1e-4), mask),
)
\end{lstlisting}

Pour des paramètres faits main (liste de dicts \texttt{\{'W', 'b'\}}, pas
Flax), le masque se construit à la main de la même forme :
\begin{lstlisting}
mask = [{'W': True, 'b': False} for _ in params]
tx = optax.masked(optax.add_decayed_weights(1e-4), mask)
\end{lstlisting}

\begin{mdframed}
\textbf{Testé} : \texttt{optax.masked(optax.add\_decayed\_weights(1e-4), mask)}
avec \texttt{mask = [\{'W': True, 'b': False\}, ...]} s'initialise et
s'applique sans erreur sur des paramètres faits main --- confirmé.
\end{mdframed}

% ═══════════════════════════════════════════════════════════════════════════
\section{Boucle d'entraînement}

Le patron standard : un \texttt{train\_step} \texttt{jit}-compilé qui prend
les paramètres et l'état de l'optimiseur, renvoie les nouveaux paramètres,
le nouvel état, et la perte.

\begin{lstlisting}
def loss_fn(params, x, y, rng, training):
    preds = model.apply({'params': params}, x, training=training, rngs={'dropout': rng})
    return jnp.mean((preds - y) ** 2)

@jax.jit
def train_step(params, opt_state, x, y, rng):
    loss, grads = jax.value_and_grad(loss_fn)(params, x, y, rng, True)
    updates, opt_state = tx.update(grads, opt_state, params)
    params = optax.apply_updates(params, updates)
    return params, opt_state, loss

@jax.jit
def eval_loss(params, x, y):
    return loss_fn(params, x, y, jax.random.key(0), False)   # dropout desactive

for step in range(n_steps):
    key, k_step = jax.random.split(key)
    params, opt_state, train_loss = train_step(params, opt_state, X_train, y_train, k_step)
\end{lstlisting}

\subsection{Full-batch vs mini-batch}
Avec un nombre de lignes informatives faible (quelques centaines à quelques
milliers, comme observé pour \texttt{lambda2p}/\texttt{lambda3p} dans ce
projet), le \textbf{full-batch} (tout le jeu d'entraînement à chaque pas) est
préférable : le bruit du mini-batch n'apporte aucun bénéfice de
généralisation à cette échelle, et le gradient exact est de toute façon peu
coûteux à calculer. Le mini-batch ne devient utile que si le jeu de données
est trop gros pour tenir en mémoire (le mur mémoire rencontré à $N=200\,000$
avec la matrice one-hot dense en est un exemple concret).

\begin{lstlisting}
def data_stream(key, X, y, batch_size):
    n = X.shape[0]
    while True:
        key, k_perm = jax.random.split(key)
        perm = jax.random.permutation(k_perm, n)
        for i in range(0, n - batch_size + 1, batch_size):
            idx = perm[i:i+batch_size]
            yield X[idx], y[idx]
\end{lstlisting}

% ═══════════════════════════════════════════════════════════════════════════
\section{Régularisation}

\subsection{Dropout}
Déjà inclus dans \texttt{ProfileMLP} ci-dessus via \texttt{nn.Dropout}. Le
point important : \texttt{deterministic=not training} --- le dropout doit
être désactivé à l'évaluation, sinon les métriques de validation sont
bruitées inutilement.

\subsection{Weight decay}
Voir \S\ref{sec:optax}. Un seul scalaire global reste une limite structurelle
similaire à celle du $\lambda$ de Ridge --- ne pas compter uniquement dessus.

\subsection{Early stopping}
Le levier le plus important pour un MLP sur un jeu de données de cette
taille : un second critère d'arrêt, indépendant du weight decay, basé sur la
perte de validation.

\begin{lstlisting}
class EarlyStopper:
    def __init__(self, patience=10, min_delta=1e-4):
        self.patience = patience
        self.min_delta = min_delta
        self.best = float('inf')
        self.counter = 0

    def step(self, val_loss) -> bool:
        """Renvoie True si l'entrainement doit s'arreter."""
        if val_loss < self.best - self.min_delta:
            self.best = val_loss
            self.counter = 0
            return False
        self.counter += 1
        return self.counter >= self.patience
\end{lstlisting}

\begin{lstlisting}
stopper = EarlyStopper(patience=30)
for step in range(max_steps):
    key, k_step = jax.random.split(key)
    params, opt_state, train_loss = train_step(params, opt_state, X_tr, y_tr, k_step)
    if step % 20 == 0:
        val_loss = float(eval_loss(params, X_va, y_va))
        if stopper.step(val_loss):
            print(f"Early stopping a l'etape {step}")
            break
\end{lstlisting}

\begin{mdframed}
\textbf{Testé} : la classe \texttt{EarlyStopper} ci-dessus, avec
\texttt{patience=3}, déclenche correctement l'arrêt sur la séquence de pertes
$[1.0,\ 0.8,\ 0.79,\ 0.79,\ 0.79,\ 0.79]$ --- à l'indice 5, une fois 3 étapes
consécutives sans amélioration $> 10^{-4}$ --- confirmé.
\end{mdframed}

% ═══════════════════════════════════════════════════════════════════════════
\section{Évaluation}

\subsection{Split train/validation}
\begin{lstlisting}
def train_val_split(X, y, val_frac=0.2, seed=0):
    n = X.shape[0]
    idx = np.random.RandomState(seed).permutation(n)
    n_val = int(n * val_frac)
    va, tr = idx[:n_val], idx[n_val:]
    return X[tr], y[tr], X[va], y[va]
\end{lstlisting}

\subsection{Corrélation de Pearson (en JAX)}
Équivalent JAX de \texttt{analysisV1.pearson}, utile si l'on veut la calculer
à l'intérieur d'une fonction \texttt{jit}-compilée.
\begin{lstlisting}
def pearson_jnp(a, b):
    a = a - jnp.mean(a)
    b = b - jnp.mean(b)
    return jnp.sum(a * b) / (jnp.sqrt(jnp.sum(a**2)) * jnp.sqrt(jnp.sum(b**2)))
\end{lstlisting}

\begin{mdframed}
\textbf{Testé} : sur $a \sim \mathcal{N}(0,1)$ et $b = 2a + 0.1\varepsilon$,
\texttt{pearson\_jnp(a,b)} $\approx 0.9999998$ --- confirmé.
\end{mdframed}

\subsection{Point d'attention : que peut-on comparer à la vérité terrain ?}
\texttt{RegressionV1.evaluate\_recovery} compare \texttt{F\_hat.ravel()} à
\texttt{protocol.F\_viab.ravel()} --- cela ne fonctionne que pour un modèle
qui produit explicitement une matrice additive par position/acide aminé
(comme Ridge). Un MLP avec une couche cachée non linéaire \textbf{n'a pas} de
telle matrice à extraire : sa sortie est une fonction non linéaire de
l'entrée one-hot, pas une somme pondérée terme à terme. Dès qu'une couche
cachée est ajoutée, la comparaison doit se faire sur les \textbf{scores
prédits} (\texttt{pearson\_jnp(v\_hat, v\_gt)}), pas sur une matrice de
poids reconstituée.

% ═══════════════════════════════════════════════════════════════════════════
\section{Exemple complet : MLP pour le modèle profile}

Cet exemple assemble tous les éléments ci-dessus sur les données réelles du
pipeline (\texttt{build\_multi\_round\_dataset}). Testé de bout en bout sur
une petite librairie synthétique ($N{=}2000$) avant rédaction : la perte
d'entraînement et de validation décroissent normalement (perte de validation
passant de $20.15$ à $1.46$ sur 300 étapes), sans erreur.

\begin{lstlisting}
import jax, jax.numpy as jnp
import numpy as np
import flax.linen as nn
import optax
import sequence_classesV1 as sc
import RegressionV1 as rg

# ---- 1. Donnees : cibles de log-enrichissement (memes que pour Ridge) ----
key = jax.random.key(11)
key, k_seq = jax.random.split(key)
sequences = jax.random.randint(k_seq, shape=(2000, 7), minval=0, maxval=20)
F_viab, F_sel, J_viab, J_sel = sc.initialize_random_weights(key=key)

protocol = sc.ProtocolV2(
    N0=1_000_000, N1=500_000, dilution_factor=10, sequences=sequences, D=100_000,
    F_viab=F_viab, J_viab=J_viab, F_sel=F_sel, J_sel=J_sel,
    noise_viab=0.1, noise_sel=0.1,
)
seqs_v, y_v, seqs_s, y_s = rg.build_multi_round_dataset(protocol, n_rounds=1)

# ---- 2. Features one-hot (profile only : L*A = 140) ----
X_v = np.eye(20, dtype=np.float32)[seqs_v].reshape(seqs_v.shape[0], -1)
y_v = np.array(y_v, dtype=np.float32)

# ---- 3. Split train/val ----
n = X_v.shape[0]
idx = np.random.RandomState(0).permutation(n)
n_tr = int(n * 0.8)
tr, va = idx[:n_tr], idx[n_tr:]
X_tr, y_tr = jnp.array(X_v[tr]), jnp.array(y_v[tr])
X_va, y_va = jnp.array(X_v[va]), jnp.array(y_v[va])

# ---- 4. Modele ----
class ProfileMLP(nn.Module):
    hidden_dim: int = 32
    @nn.compact
    def __call__(self, x, training: bool):
        x = nn.Dense(self.hidden_dim)(x)
        x = nn.relu(x)
        x = nn.Dropout(rate=0.1, deterministic=not training)(x)
        x = nn.Dense(1)(x)
        return x.squeeze(-1)

model = ProfileMLP()
key, k_init, k_drop = jax.random.split(key, 3)
params = model.init({'params': k_init, 'dropout': k_drop}, X_tr[:1], training=False)['params']

# ---- 5. Perte, optimiseur, schedule ----
def loss_fn(params, x, y, rng, training):
    preds = model.apply({'params': params}, x, training=training, rngs={'dropout': rng})
    return jnp.mean((preds - y) ** 2)

schedule = optax.warmup_cosine_decay_schedule(0.0, 1e-3, warmup_steps=20, decay_steps=300)
tx = optax.adamw(learning_rate=schedule, weight_decay=1e-3)
opt_state = tx.init(params)

@jax.jit
def train_step(params, opt_state, rng):
    loss, grads = jax.value_and_grad(loss_fn)(params, X_tr, y_tr, rng, True)
    updates, opt_state = tx.update(grads, opt_state, params)
    params = optax.apply_updates(params, updates)
    return params, opt_state, loss

@jax.jit
def eval_loss(params):
    return loss_fn(params, X_va, y_va, jax.random.key(0), False)

# ---- 6. Boucle avec early stopping ----
stopper = EarlyStopper(patience=30)
for step in range(300):
    key, k_step = jax.random.split(key)
    params, opt_state, train_loss = train_step(params, opt_state, k_step)
    if step % 20 == 0:
        val_loss = float(eval_loss(params))
        print(f"step {step:3d}  train {float(train_loss):.4f}  val {val_loss:.4f}")
        if stopper.step(val_loss):
            print(f"Early stopping a l'etape {step}")
            break
\end{lstlisting}

% ═══════════════════════════════════════════════════════════════════════════
\section{Pièges et bonnes pratiques spécifiques à ce projet}

\begin{itemize}
    \item \textbf{Clés aléatoires} : ne jamais réutiliser \texttt{key} sans
          \texttt{split} --- chaque appel à \texttt{model.init}, chaque pas
          d'entraînement, chaque \texttt{jax.random.permutation} doit recevoir
          une sous-clé fraîche. C'est la même discipline que
          \texttt{Protocol.\_next\_key()} dans \texttt{sequence\_classesV1.py}.
    \item \textbf{float32 par défaut} : JAX désactive le float64 par défaut
          (contrairement à NumPy). Les tableaux \texttt{np.float64} passés à
          JAX sont silencieusement tronqués en float32 --- généralement sans
          conséquence ici, mais à garder en tête si des NaN apparaissent après
          beaucoup d'étapes d'entraînement.
    \item \textbf{Dropout et \texttt{rngs}} : oublier de passer
          \texttt{rngs=\{'dropout': ...\}} à \texttt{model.apply} lève une
          erreur explicite si \texttt{training=True} --- ce n'est pas un bug
          silencieux, contrairement à d'autres frameworks.
    \item \textbf{Évaluation sans dropout} : toujours appeler
          \texttt{model.apply(..., training=False)} pour l'évaluation/la
          validation, sinon la perte de validation est artificiellement
          bruitée par le dropout actif.
    \item \textbf{Vérifier le modèle linéaire avant le MLP} : entraîner
          d'abord un \texttt{ProfileMLP} \textbf{sans} couche cachée (un seul
          \texttt{nn.Dense} direct de 140 vers 1) avec AdamW, et vérifier que
          le score de corrélation obtenu est proche de celui de Ridge au même
          niveau de régularisation. Cela valide la boucle d'entraînement et le
          pipeline de données \emph{avant} d'introduire la non-linéarité --- si
          quelque chose ne va pas ensuite, on sait que c'est la couche cachée,
          pas la plomberie.
    \item \textbf{\texttt{jit} et formes variables} : si la taille du batch
          change d'un appel à l'autre (par exemple en fin d'époque avec un
          reste plus petit), JAX recompile --- coûteux si répété souvent.
          Préférer une taille de batch fixe (tronquer le reste) ou le
          full-batch (\S 8.1).
\end{itemize}

% ═══════════════════════════════════════════════════════════════════════════
\section{Aide-mémoire récapitulatif}

\begin{center}
\begin{tabular}{p{5cm}p{9.5cm}}
\toprule
\textbf{Fonction} & \textbf{Rôle} \\
\midrule
\texttt{jax.random.split(key, n)} & Dérive $n$ sous-clés indépendantes \\
\texttt{jax.grad} / \texttt{value\_and\_grad} & Gradient (et valeur) d'une fonction scalaire \\
\texttt{jax.jit} & Compile XLA d'une fonction \\
\texttt{jax.vmap} & Vectorise une fonction écrite pour un seul exemple \\
\texttt{jax.tree\_util.tree\_map} & Applique une fonction à toutes les feuilles d'un pytree \\
\texttt{jax.nn.relu / gelu / sigmoid / softmax} & Activations \\
\texttt{flax.linen.Dense} & Couche dense (linéaire + biais) \\
\texttt{flax.linen.Dropout} & Dropout, contrôlé par \texttt{deterministic} \\
\texttt{model.init(rngs, x, ...)} & Initialise les paramètres \\
\texttt{model.apply(\{'params':...\}, x, ...)} & Passe avant \\
\texttt{optax.adamw} & Optimiseur Adam + weight decay découplé \\
\texttt{optax.exponential\_decay} & Schedule de taux d'apprentissage exponentiel \\
\texttt{optax.warmup\_cosine\_decay\_schedule} & Schedule warmup + cosinus \\
\texttt{optax.clip\_by\_global\_norm} & Écrêtage de gradient \\
\texttt{optax.chain} & Compose plusieurs transformations de gradient \\
\texttt{optax.masked} & Applique une transformation seulement sur certains paramètres \\
\texttt{optax.add\_decayed\_weights} & Weight decay pur (utilisé avec \texttt{masked}) \\
\texttt{tx.init(params)} / \texttt{tx.update(grads, opt\_state, params)} & Cycle de vie de l'optimiseur \\
\texttt{optax.apply\_updates(params, updates)} & Applique les mises à jour aux paramètres \\
\bottomrule
\end{tabular}
\end{center}

\begin{mdframed}
\textbf{Résumé} : chaque exemple de code de ce document a été exécuté et
vérifié avant rédaction (JAX 0.6.2, Optax 0.2.8, Flax 0.10.7). Le fil
conducteur recommandé : partir du modèle linéaire pur (une seule couche
\texttt{Dense}, sans activation cachée) pour valider la boucle
d'entraînement contre Ridge, puis n'introduire la non-linéarité, le dropout
et l'early stopping qu'une fois cette base confirmée correcte.
\end{mdframed}

\end{document}
